% Full constructive GPSS bounded/ray complement proofs and endpoint correction.

\subsection{The GPSS complement construction and its exact endpoint domain}
\label{sec:annex-p09-gpss-decomposition-complement}

\begin{annexitem}{annex:p9.gpss.finite.bounded.interval}{thm:p09-gpss-finite-bounded-interval}{Successors, completeness and finite bounded intervals}
\begin{proof}
If \(a=b\), the closed interval is a singleton. Suppose
\(a<b\). Start with \(a_0=a\). Whenever \(a_n<b\),
the point \(a_n\) is not maximal and therefore has an
immediate successor \(a_{n+1}\). That successor is at
most \(b\); otherwise \(b\) would lie strictly between
the two adjacent points. If this process reaches \(b\),
its finite successor list exhausts \([a,b]\), since
there is no point between consecutive members.

Suppose it never reaches \(b\). The resulting strictly
increasing sequence is bounded below by \(a\) and above
by \(b\), so completeness gives its least upper bound
\(s\in T\). No sequence member equals \(s\), since
it has a larger next member. Also \(s>a\), so discreteness
gives an immediate predecessor \(p<s\). The point
\(p\) cannot be an upper bound of the sequence, by
leastness of \(s\). Some \(a_n\) therefore satisfies
\(p<a_n<s\), contrary to adjacency. This contradiction
proves termination of the successor process and finiteness
of the interval. The first and last member of each nonempty
subset of a finite interval exist. The open interval of
adjacent points has no member, as claimed.
\end{proof}
\end{annexitem}

\begin{annexitem}{annex:p9.gpss.canonical.prefix}{thm:p09-gpss-canonical-prefix}{The future-labelled earliest feasible chain}
\begin{proof}
The temporal recursion for \(\Lambda_i\) expresses
existence of the remaining positive chain after the current
point, without any final gap constraint or bound \(b\).
This follows backwards from \(\Lambda_m=\top\): each
strict until chooses the next labelled point and enforces
the indicated preceding gap. All these formulas remain
future-only and have no endpoint parameter.

Assume first that a satisfying bounded chain
\(v_1<\cdots<v_m\) exists. Each \(v_i\) satisfies
\(\Theta_i\land\Lambda_i\), because the rest of that
same chain witnesses the remaining future formula. The
first selector \(\mu_1\) is consequently defined and
is at most \(v_1\). Inductively, if
\(\mu_{i-1}\leq v_{i-1}\), then \(v_i\) is a
qualifying point strictly between \(\mu_{i-1}\) and
\(b\). Thus \(\mu_i\) is defined and
\(\mu_i\leq v_i\). This proves the coordinatewise
inequalities. The first required gap \((a,v_1)\) is
labelled \(\Phi_0\), so its subinterval
\((a,\mu_1)\) has no \(\neg\Phi_0\)-point.
The last-failure selector in \(V_m\) is undefined.
All the Boolean conditions of \(V_m\) follow.

Conversely suppose \(V_m\) holds, with \(m>0\).
All selected points exist and are strictly increasing
inside \((a,b)\). Their point labels hold by the
selector definitions. The first gap has its required
label by absence of the last \(\neg\Phi_0\) event.
For \(i\geq2\), the preceding point \(\mu_{i-1}\)
satisfies \(\Lambda_{i-1}\). It therefore has some
later point \(w\) satisfying
\(\Theta_i\land\Lambda_i\), with \(\Phi_{i-1}\)
throughout \((\mu_{i-1},w)\). If \(w<b\),
the earliest bounded choice \(\mu_i\) is at most
\(w\). If \(w\geq b\), the same inequality
holds because \(\mu_i<b\). Hence
\((\mu_{i-1},\mu_i)\) is contained in that valid
gap and satisfies \(\Phi_{i-1}\). This proves
every intervening gap, even if the future viability
witness originally lay beyond the upper bound.
The selected chain is a valid bounded chain.

For \(m=0\), the chain has no required point
or preceding gap, so its existence assertion is
true on the given admissible interval. This agrees
with the defined constant test. The preceding
inequalities prove minimality whenever the chain
is nonempty. Every first selector has a least
qualifying point when one exists because its
search interval is bounded and finite. Empty
intervals or missing candidates give the specified
undefined statuses rather than arbitrary witnesses.
\end{proof}
\end{annexitem}

\begin{annexitem}{annex:p9.gpss.determining.events}{thm:p09-gpss-determining-events}{The full last-gap-failure recursion}
\begin{proof}
Induct on the number \(r\) of internal points.
For rank zero, the decomposition is exactly the
universal condition \(\Phi_0\) on its open
interval. In a finite interval this holds if
and only if there is no last \(\neg\Phi_0\)
occurrence, proving the base decision formula.

Suppose \(r>0\), and let \(\ell\) be the
last \(\neg\Phi_r\)-point in \((a,b)\),
if one exists. If it does not, the final gap
condition holds everywhere in the interval.
The decomposition is then equivalent to
existence of its entire positive chain with
that final gap omitted, exactly \(V^0(a,b)\)
by the preceding lemma. This proves the first
branch of the decision expression.

Now let \(\ell\) be defined. Every valid
decomposition chain must finish at or after
\(\ell\). Otherwise this false final-gap
point would lie strictly between its last
witness and \(b\). Conversely, once a chain
finishes at or after \(\ell\), its final
gap is automatically valid, because there
is no later \(\neg\Phi_r\)-point before
\(b\).

Take a valid chain. If \(\ell\) lies
strictly between its \(j\)-th and
\((j+1)\)-st points, with the zeroth
point understood as \(a\), then
\(0\leq j<r\). The initial part gives
\(D_{W_j}(a,\ell)\), the cut point
itself satisfies \(\Phi_j\), and the
remaining chain after \(\ell\) gives
\(V^{j}(\ell,b)\). The inductive
decision for \(W_j\) and the defined
test at \(\ell\) therefore satisfy
the corresponding gap branch. The
\(\Phi_j(\ell)\) test is necessary
because the original open gap includes
this inserted point.

If instead \(\ell\) equals the
\(j\)-th witness, where \(1\leq j\leq r\),
the part strictly before it gives
\(D_{W_{j-1}}(a,\ell)\), the point
label is \(\Theta_j(\ell)\), and
the rest gives \(V^{j}(\ell,b)\).
This is exactly the corresponding
point branch. For \(j=r\), the
last test is the empty positive-chain
test and imposes no further witness.
No \(\Phi_j\)-condition is added
at a genuine point witness, because
the adjacent gap universals exclude
their endpoints. These alternatives
exhaust all positions of \(\ell\)
relative to a chain finishing at
or after it.

For the reverse implication, a true
gap branch supplies its valid initial
decomposition by induction, its
\(\Phi_j\)-labelled cut point,
and a valid remaining positive chain
by the canonical-chain lemma. Join
the two chains. The two pieces of
the split \(\Phi_j\)-gap and the
label at the cut together enforce
the original whole gap. All other
point and gap conditions are retained.
The resulting last witness is after
\(\ell\), so its final gap is
valid. A true point branch supplies
the valid initial decomposition,
inserts \(\ell\) as its next
labelled witness, and appends the
remaining positive chain. The final
witness is at or after \(\ell\),
so this also satisfies the entire
original decomposition. Both
directions of the recursive decision
are proved.

All recursive calls use ranks
\(j<r\). Each level adds only
finitely many last selectors, tests
and finite canonical chains. Thus
induction gives a finite list and
an explicit terminating procedure
for constructing it. With endpoint
subterms shared, a bound on the
number of new event nodes can be
defined by \(B(0)=1\) and
\[
                  B(r)=(r+3)^2+\sum_{j<r}B(j)
                  \qquad(r\geq1).
\]
Indeed the \(r\) prefix decisions
can be reused in the two sums;
the last event, two sets of \(r\)
tests, the full positive-chain test
and all suffix positive-chain tests
use at most \((r+3)^2\) further
nodes. This bounds event nodes,
not the length of copied label
formulas. Those are also finite,
being explicit until recursions.

Every test is applied to a selected
internal point. No condition has
been put on the external endpoints'
predicate values. If an event bound
in an inactive recursive branch
is undefined, its descendants are
undefined by the stipulated semantics;
the surrounding definedness guard
makes that branch irrelevant. The
root equivalence is asserted only
on its actual admissible bounds.
Finally, the decision uses solely
Boolean operations on definedness
bits. Its truth is consequently
determined by the complete statuses
and order of the finite list,
including every equality among
defined event values.
\end{proof}
\end{annexitem}

\begin{annexitem}{annex:p9.gpss.bounded.complement}{thm:p09-gpss-bounded-complement}{Compilation of every event status and weak-order case}
\begin{proof}
Form the finite event list closed
under subterms from the preceding
construction. Choose which of
its nonroot terms are defined,
and choose an ordered partition
of the roots and defined terms.
The roots must be defined and
ordered \(x<y\). Every defined
nonroot value must be strictly
between them. A block in the
partition denotes equality, and
the block order denotes strict
order. There are finitely many
such status/type choices.

We compile one choice. Work
through its term graph in
construction order. For a first
or last selector, if either bound
is declared undefined, its output
must be undefined; reject a case
which declares otherwise. If both
bounds are defined but not strictly
ordered, the same rule applies.
In these situations no occurrence
condition is imposed. If the
bounds are defined and ordered,
there are two remaining cases.

For a defined first selector with
declared value \(v\), require
its lower bound to be below
\(v\), and \(v\) below its
upper bound. Put the selector
label at \(v\), and require
its negation throughout the
open interval from the lower
bound to \(v\). For a defined
last selector, use the same
point label and require its
negation from \(v\) to the
upper bound. Reject inconsistent
order requirements. For an
undefined selector with defined
ordered bounds, require its
negated label throughout the
whole open interval between
those bounds. This last clause
is exact because a nonempty
occurrence set in that finite
interval necessarily has both
extrema. No nonattainment case
is being suppressed under the
current model hypotheses.

For a test node, an undefined
argument forces an undefined
result. With a defined argument,
a defined result must lie in
the same equality block and
its label must be true at
that point; otherwise reject
the order case. An undefined
test with a defined argument
requires the label's negation
at the argument point. These
are all its semantic cases.
Test arguments are internal
terms, so no point label has
been required at \(x\) or
\(y\).

After these finite checks the
case consists only of its
weak order, finitely many
future-label point conditions,
and finitely many universal
future-label conditions on
open intervals whose bounds
are among the declared values.
For each internal equality
block introduce one fresh
point variable. Order these
variables strictly between
\(x,y\). The original term
names in one block are
aliases for that variable,
so they impose no further
quantified points.

Conjoin at each point all
its required point labels.
For every universal condition
on \((s,t)\), conjoin its
label at each internal block
strictly between those bounds,
and on each merged open gap
between consecutive blocks
contained in \((s,t)\).
Do not add it at a block
equal to \(s\) or \(t\).
Every remaining empty label
list is true. These are
finite conjunctions of labels
or their negations, all in
the exact permitted \(L(U)\)
class. Labels are copied with
fresh bound variables, so
only their free current point
is replaced by the point or
gap variable.

This is one GPSS decomposition.
It includes the guard \(x<y\),
has only internal labelled
points and labelled open
gaps, and has no additional
free parameter or event
symbol. A label conjunction
which is inconsistent merely
makes that decomposition
unsatisfiable; detecting all
such inconsistencies is not
needed by the procedure.

To verify its meaning, first
take actual event values
realizing the case. First
and last occurrences satisfy
their point and absence
conditions. Undefined
selectors with good bounds
have empty occurrence sets,
and tests satisfy exactly
their specified truth cases.
The universal intervals
are partitioned into the
merged gaps and internal
points that they cover.
Thus the corresponding
merged points witness the
compiled decomposition.

Conversely suppose a compiled
decomposition has witnesses.
Assign each defined term
the point of its block and
each other term the undefined
value. Induct along the term
graph. Roots have their fixed
values. A selector with failed
bounds is undefined as declared.
For good ordered bounds, the
point label and absence before
or after a declared value make
it the actual first or last
occurrence. The absence on the
whole interval in an undefined
case means there is no such
occurrence. Test nodes follow
their point label and equality
constraints. Each universal
constraint follows from its
labels on all the constituent
open gaps and inserted points.
Hence these are exactly the
actual event statuses and
values, with the declared
weak order.

Every assignment with \(x<y\)
has one and only one actual
status/type case. Evaluate
the finite Boolean expression
\(T_W\) on the case's
definedness bits. Retain the
compiled decomposition exactly
when that Boolean value is
false. The determining lemma
then proves that the disjunction
of retained cases is equivalent
to \(\neg D_W\) on the
admissible interval. Outside
it, all output decompositions
and the guarded target are
false. This proves the
displayed relative equivalence.

An empty retained list is
false. If the actual open
interval is empty, all strict
selectors are undefined; the
same event rules still apply,
with vacuous universal gaps.
Rank zero and positive rank
therefore receive their exact
opposite complement values
there. The source's finite
event-order argument is thus
supplied with an explicit
grammar, complete partial
semantics, a terminating
determining-list construction
and a back-translation for
every case
\cite[Proposition~2.4 outline,
       p.~171]{GabbayPnueliShelahStavi1980}.

For \cref{def:p09-gpss-last-failure-example},
the last failed final-gap point
is two. If the point label
holds there, the rank-one
equality branch uses it as
the sole witness; the final
open gap excludes it and is
valid. If the label holds
only at one, that earlier
candidate leaves the false
point two inside its final
gap, while no candidate at
or after two exists. The
decision is false in that
case. On adjacent endpoints
no positive-rank witness
chain exists at all, so its
relative complement is the
full admissible interval
case. These verify the
worked equality and vacuity
boundaries directly.
\end{proof}
\end{annexitem}

\begin{annexitem}{annex:p9.gpss.suffix.complement}{thm:p09-gpss-suffix-complement}{The strict next-point unfolding for a future ray}
\begin{proof}
The proved suffix translation has
the form \(B=H_0\) conjoined
with finitely many \(F\Psi\)
and \(GF\Psi\) factors.
Its outer Boolean skeleton has
only strict modalities and
truth constants as leaves:
\(H_0\) is an until formula
or a strict \(G\)-formula.
In particular it has no exposed
current-point atom. The same
is true of \(E=\neg B\).
Keep the designated Boolean
constants until the final
expansion into \(L(U)\).

If \(x\) has no later point,
every \(F\)-leaf and until
leaf is false and every
\(G\)-leaf is true. Evaluating
the outer Boolean skeleton
therefore gives exactly
\(e_{\mathrm{last}}\), independently
of the input value at \(x\).

If \(x\) is not last,
discreteness gives its immediate
successor \(s\). The later
points of \(x\) consist of
\(s\) and the points strictly
later than \(s\). This gives
\(FA\) at \(x\) exactly
when \(A\lor FA\) holds
at \(s\), and \(GA\)
at \(x\) exactly when
\(A\land GA\) holds there.
For a strict until, its
witness is either \(s\),
requiring only \(C(s)\),
or is later than \(s\).
The latter case requires
\(A(s)\) and precisely the
same until formula at \(s\).
This proves the third unfolding
clause with its exact operand
order. Boolean recursion then
gives
\[
                         E(x)\quad\Longleftrightarrow\quad L_E(s)
\]
for every nonlast \(x\).
All formulas in \(L_E\) are
future-only and have no reference
to the former point \(x\).

The one-point output decomposition
chooses \(s>x\) with no point
strictly between \(x,s\),
because its first gap label is
false. Thus its witness is
exactly the immediate successor,
and its point label tests
\(L_E(s)\). Its true final
gap adds no condition. At a
last point this decomposition
is false. The optional zero-point
decomposition with gap false
is true exactly at last points.
Adding it precisely when
\(e_{\mathrm{last}}\) is true
therefore gives \(E=\neg B\)
at every point. By the previous
suffix equivalence this is
\(\neg D(x,+\infty)\).

The global cofinal conjuncts
are included in \(B\) by
their exact \(F\Psi\land GF\Psi\)
translation. Hence empty or
nonempty cofinal lists require
no extra implicit endpoint
assumption. This construction
avoids unbounded last-occurrence
terms entirely. It is a
proved return through the
source's future translation
and remains in the exact
decomposition grammar, rather
than enlarging its label class.
\end{proof}
\end{annexitem}

\begin{annexitem}{annex:p9.gpss.prefix.complement.obstruction}{thm:p09-gpss-prefix-complement-obstruction}{The precise obstruction at a formal left endpoint}
\begin{proof}
Work on the fixed discrete complete
order \(\mathbb Z\), with upper
endpoint zero for the prefix case.
For each binary sequence
\(w=(w_1,w_2,\ldots)\), interpret
\(q_0(-n)\) by \(w_n\),
interpret \(q_0(k)\) as true
for every \(k\geq0\), and
fix every other predicate to be
false everywhere. These restrictions
are legitimate instances of the
arbitrary valuations in any alleged
uniform equivalence.

A cylinder is the set of all binary
sequences extending one finite
initial bit string. A set is closed
if every sequence outside it has
a cylinder neighborhood disjoint
from it. Each future-only label
evaluated at a fixed integer \(k\)
depends on only finitely many bits:
for \(k<0\), the varying part
of its future tail consists of
the positions \(k,\ldots,-1\);
for \(k\geq0\), the tail is
completely fixed. The future-locality
lemma therefore makes that label's
truth set a finite union of
cylinders. Both it and its
complement are closed.

Fix now a particular finite tuple
of witnesses for any prefix
decomposition. Its point conditions
and all its universal gap conditions
form an intersection of closed
sets: each universal is an
intersection of the label conditions
at the relevant fixed integers.
An arbitrary intersection of closed
sets is closed, since failure of
one condition already gives a
disjoint cylinder neighborhood.
There are countably many finite
integer witness tuples. Thus the
truth set of that decomposition
on the chosen valuation space
is a countable union of closed
sets. The same holds for a
finite disjunction of such
decompositions.

For whole-order decompositions the
same argument applies to all
their gaps, including the final
unbounded gap. Each allowed
future-cofinal condition is constant
on this valuation space: evaluate
its equivalent \(F\Psi\land GF\Psi\)
at zero, whose entire labelled
future tail is fixed. Adding
such constants preserves closedness
for a fixed witness tuple.
Consequently a finite disjunction
of whole-order decompositions also
gives a countable union of closed
sets. This reasoning permits any
finite future formulas as labels,
not only atomic ones.

The input decomposition
\(\exists z<0\,\forall u<z\,q_0(u)\)
holds exactly when \(w\) has
only finitely many zeros. If
there are finitely many, choose
a negative \(z\) below all
their positions. Conversely a
valid \(z\) leaves zeros only
at finitely many integers between
\(z\) and minus one. The
whole-order version with unrestricted
\(z\) has the same condition:
finitely many zeros permit the
same negative witness, while
infinitely many zeros occur below
every integer and prevent any
witness. Their complements on
this valuation space therefore
give the set \(\mathcal I\)
of sequences with infinitely
many zeros.

We show directly that
\(\mathcal I\) is not a
countable union of closed sets.
Suppose \(\mathcal I=\bigcup_n C_n\)
with each \(C_n\) closed.
Build nested finite strings.
Given the current string,
extend it by ones forever,
obtaining a sequence with
only finitely many zeros.
It is outside \(C_n\).
Closedness gives a sufficiently
long finite prefix of this
sequence whose cylinder avoids
\(C_n\). Extend that prefix
by one zero to form the
next string. The strings
remain compatible, their
lengths strictly increase,
and they add a new zero
at each step. Their union
is therefore an infinite
sequence in \(\mathcal I\).
It avoids every \(C_n\),
because its later prefixes
stay in the cylinder chosen
to avoid that set. This
contradiction proves the
claim without an unproved
topological separation assertion.

An alleged finite decomposition
disjunction for either complement
would make \(\mathcal I\)
a countable union of closed
sets by the earlier paragraphs.
That is impossible. Both
counterexamples are legal
members of the original
future-label class: one
point labelled true, first
gap \(q_0\), final gap true,
and no cofinal requirement.
Thus this is an actual
failure of the broader
formal-left-endpoint closure
claim even on a discrete
complete order. A dual past
label class or coinitial
extension is not silently
substituted for it here.
\end{proof}
\end{annexitem}

\begin{annexitem}{annex:p9.gpss.discrete.future.completeness}{thm:p09-gpss-discrete-future-completeness}{The full future regional induction and the point-zero corollary}
\begin{proof}
Induct simultaneously on the GPSS
regional stages with actual lower
endpoint and either actual or
formal upper endpoint. Stage zero
is empty. For a positive
constructor, the induction hypotheses
supply finite decomposition lists
for all its regional children.
Their lower endpoints are still
actual points: the original
lower endpoint or the newly
quantified witness. Apply the
already proved finite positive
intersection and concatenation
construction. Point literals
are permitted future labels,
and empty child lists and
empty disjunctions have their
proved unit and zero meanings.
This gives a finite list for
the positive component.

A negative component is the
negation of this finite
positive disjunction. Complement
each of its decompositions
using the bounded or future-ray
theorem according to its upper
endpoint, and conjoin the
resulting finite lists by
positive closure. All these
operations are on the same
admissible interval. The empty
positive disjunction has the
true one-gap complement, and
an empty factor in a conjunction
gives false. Thus all signed
constructor cases and all
finite stages are covered.
This proves the stated Main
Lemma on the interval types
actually needed for the
future construction.

For a source future first-order
formula \(\varphi(x)\), use
\cref{thm:p09-gpss-future-special-form}
to obtain a finite disjunction
of conjunctions of point literals
at \(x\) and after components
at \(x\). Replace each after
component by its just-proved
finite decomposition list,
then by its exact future
expression from
\cref{thm:p09-gpss-suffix-translation}.
Point literals translate to
the corresponding temporal
atoms or their negations.
Finite Boolean operations
remain in \(L(U)\). All
equivalences hold at every
point of every discrete
complete model, including
its endpoints. This proves
the full Theorem~2.2 target.

Finally let \(\varphi(x)\)
be an arbitrary first-order
formula. Rename its bound
variables fresh and restrict
each quantifier to values
at or after \(x\), obtaining
a source future formula.
At point zero of any
\(\omega\)-model, every
domain point already satisfies
that restriction. Induction
through the formula shows
that its value there is
unchanged. Applying the
future result gives the
required temporal formula
at zero, which is exactly
Theorem~2.1. No equivalence
to the original unrestricted
formula at other points is
asserted.

The finite bounded-interval
lemma and the successor
unfolding are the discrete
steps in this proof. They
do not furnish first/last
attainment or a next point
on an arbitrary complete
order. The broader complete
and gap cases therefore
remain separate source
constructions
\cite[pp.~170--171]{GabbayPnueliShelahStavi1980}.
\end{proof}
\end{annexitem}
